\section{The Page That Costs Four Statements}
\label{sec:ch10-the-page-that-costs-four}

\index{N+1@N+1!at the entity route|(}%
Here is the request chapter~\ref{ch:data-without-n-plus-1} has been sending
since the graph was one service, going through the router at the graph
chapter~\ref{ch:second-third-subgraph} finished building:

\begin{minted}{text}
POST /graphql HTTP/1.1
Host: localhost:3002
Content-Type: application/json

{"query":"{ sessions(first: 10) { nodes { title speaker { name } } } }"}
\end{minted}

\begin{minted}[fontsize=\footnotesize]{json}
{"data":{"sessions":{"nodes":[
  {"title":"Schemas That Outlive Their Authors",
   "speaker":{"name":"Ada Fischer"}},
  {"title":"Reading a Query Plan",
   "speaker":{"name":"Ada Fischer"}},
  {"title":"Paging Without Offsets",
   "speaker":{"name":"Bruno Kaminski"}},
  {"title":"The Bank That Split Its Graph",
   "speaker":{"name":"Chidi Okafor"}}]}}}
\end{minted}

\index{statement count!the seam costs two more than it should}%
Correct, complete, HTTP 200, and no error key.
Chapter~\ref{ch:second-third-subgraph} measured this exact request at two
statements, one in each of two services, and made a point of the two being
chapter~\ref{ch:data-without-n-plus-1}'s number surviving a seam. Change one
parameter and the body beneath it, in one method, and it is four:

\begin{center}
\begin{tabular}{lrr}
\toprule
Service & Chapter~\ref{ch:second-third-subgraph} & Now \\
\midrule
Sessions, on 5001 & 1 & 1 \\
Speakers, on 5002 & 1 & 3 \\
Ratings, on 5003 & 0 & 0 \\
\midrule
Total & 2 & 4 \\
\bottomrule
\end{tabular}
\end{center}

\index{reference resolver!the naive form}%
What changed is the DataLoader behind the Speakers service's reference
resolver. Chapter~\ref{ch:first-subgraph} put one there when it wrote this
book's first entity route, and measured what it was worth against three
representations sent by hand;
chapter~\ref{ch:second-third-subgraph} moved the file to the service that owns
the rows now and left it as it was. This is the version without it, which is
what a reference resolver looks like when you write the obvious thing:

\begin{minted}[highlightlines={5,17,18,19}]{csharp}
    [GraphQLIgnore]
    public static async Task<Speaker?> ResolveSpeakerReferenceAsync(
        string id,
        INodeIdSerializer serializer,
        SpeakerContext db,
        CancellationToken ct)
    {
        var nodeId = serializer.Parse(id, typeof(int));

        if (nodeId.TypeName != nameof(Speaker))
        {
            return null;
        }

        var key = (int)nodeId.InternalId!;

        return await db.Speakers
            .Where(s => s.Id == key)
            .FirstOrDefaultAsync(ct);
    }
\end{minted}

\index{reference resolver!decodes a key and fetches a row}%
Nothing about that is careless. It decodes the key the way
chapter~\ref{ch:first-subgraph} established, it checks the type name inside the
key before spending the integer, and it fetches one row for one key, which is
what the method signature says it is for. Every resolver in
chapter~\ref{ch:hot-chocolate-zero} was written like this and every one of them
was right to be.

\subsection{What Else Changed}

\index{N+1@N+1!leaves no trace outside the service}%
Nothing else changed. That is the part worth sitting with, because it is what
makes this failure different from every other one this book has shown you.

\index{schema!the omission does not reach it}%
The Speakers service exports the same schema, byte for byte. The graph
composes, with no message. The router serves the same types and the same root
fields. The query plan is the same \code{Sequence} of two steps with the same
\code{BatchEntity} fetch naming the same field path. The router makes the same
number of HTTP requests it made before, which is one. And the response, as you
have just seen, is identical.

\index{observability!the router cannot see this}%
So there is no vantage point outside the Speakers service from which this is
visible at all. From where the router stands the graph answered one query with
two subgraph calls, and that is true. The only instrument that can see the rest
is the one inside the process doing the work, and in this book that is
chapter~\ref{ch:data-without-n-plus-1}'s statement log.
Chapter~\ref{ch:blame-routing} is the chapter about finding out which subgraph
is slow, and this is the failure mode it will be hunting: a service that is
correct, is not erroring, is not timing out, and is doing four times the work
somebody thinks it is.

\index{N+1@N+1!the count is somebody else's page size}%
One more property before the mechanism, because it changes how alarming the
four is. Three of those statements are one per speaker on the page, and the
page is the Sessions service's. Nothing in the Speakers service chose it, and
nothing in the Speakers service can see it: this service was handed a list and
the list is as long as another team's default page size, or as long as whatever
a client asked that team for.
